\documentclass[aspectratio=169,10pt]{beamer}
% ---------------------------------------------------------------
% R4.02 Qualite de developpement -- preambule commun des diapositives
% Charte reprise de R3.07 (metropolis, dark tech)
% Palette IUT partagee (violet du logo + accent or)
% Utilisation :
%   \documentclass[aspectratio=169,10pt]{beamer}
%   % ---------------------------------------------------------------
% R4.02 Qualite de developpement -- preambule commun des diapositives
% Charte reprise de R3.07 (metropolis, dark tech)
% Palette IUT partagee (violet du logo + accent or)
% Utilisation :
%   \documentclass[aspectratio=169,10pt]{beamer}
%   % ---------------------------------------------------------------
% R4.02 Qualite de developpement -- preambule commun des diapositives
% Charte reprise de R3.07 (metropolis, dark tech)
% Palette IUT partagee (violet du logo + accent or)
% Utilisation :
%   \documentclass[aspectratio=169,10pt]{beamer}
%   \input{../../latex/common/r402-slides-preamble.tex}
%   \graphicspath{{../../latex/common/assets/}{img/}{figs/}}
% ---------------------------------------------------------------
\usetheme{metropolis}
\usepackage{appendixnumberbeamer}

\usepackage[T1]{fontenc}
\usepackage[utf8]{inputenc}
\usepackage[french]{babel}
\usepackage[scale=0.94]{tgheros}
\renewcommand{\familydefault}{\sfdefault}
\usepackage{amsmath,amssymb,amsfonts}
\usepackage{graphicx}
\usepackage{booktabs}
\usepackage{multirow}
\usepackage{tikz}
\usetikzlibrary{shapes.geometric, shapes, arrows, arrows.meta, positioning, calc, fit}
\usepackage{xcolor}
\usepackage{listings}
\usepackage[skins,listings]{tcolorbox}
\usepackage{qrcode}
\usepackage{eurosym}
\usepackage{hyperref}

% ---------------------------------------------------------------
% PALETTE (dark-mode tech: near-black canvas, violet (logo IUT) + gold accents)
% ---------------------------------------------------------------
\definecolor{qabg}{HTML}{0E1016}      % canvas background
\definecolor{qapanel}{HTML}{181B24}   % card / block panel
\definecolor{qapanelalt}{HTML}{20242F}% secondary panel tint
\definecolor{qaline}{HTML}{2B303D}    % hairline rules / table rules
\definecolor{qatext}{HTML}{ECEEF2}    % primary text (off-white)
\definecolor{qamuted}{HTML}{8B93A6}   % secondary / caption text
\definecolor{qaaccent}{HTML}{D93ED7}  % primary accent -- violet du logo IUT, eclairci pour fond sombre
\definecolor{qaaccent2}{HTML}{F2B84B} % secondary accent -- or/ambre
\definecolor{qacodekw}{HTML}{C792EA}  % code : mots-cles
\definecolor{qacodestr}{HTML}{C3E88D} % code : chaines
\definecolor{qacodecom}{HTML}{7F8AA3} % code : commentaires
\definecolor{qacodeann}{HTML}{F2B84B} % code : annotations / javadoc

\setbeamercolor{background canvas}{bg=qabg}
\setbeamercolor{normal text}{fg=qatext}
\setbeamercolor{palette primary}{bg=qabg, fg=qamuted}
\setbeamercolor{frametitle}{bg=qabg, fg=qatext}
\setbeamercolor{progress bar}{fg=qaaccent, bg=qaline}
\setbeamercolor{alerted text}{fg=qaaccent2}
\setbeamercolor{block title}{bg=qapanel, fg=qaaccent}
\setbeamercolor{block body}{bg=qapanel, fg=qatext}
\setbeamercolor{block title alerted}{bg=qapanel, fg=qaaccent2}
\setbeamercolor{block body alerted}{bg=qapanel, fg=qatext}
\setbeamercolor{block title example}{bg=qapanelalt, fg=qaaccent}
\setbeamercolor{block body example}{bg=qapanelalt, fg=qatext}
\setbeamercolor{example text}{fg=qaaccent}
\setbeamercolor{itemize item}{fg=qaaccent}
\setbeamercolor{itemize subitem}{fg=qaaccent2}
\setbeamercolor{enumerate item}{fg=qaaccent}
\setbeamercolor{enumerate subitem}{fg=qaaccent2}
\setbeamercolor{caption}{fg=qamuted}
\setbeamercolor{page number in head/foot}{fg=qamuted, bg=qabg}
\setbeamercolor{title}{fg=qatext}
\setbeamercolor{subtitle}{fg=qamuted}
\setbeamercolor{author}{fg=qatext}
\setbeamercolor{institute}{fg=qamuted}
\setbeamercolor{date}{fg=qamuted}
\setbeamercolor{section in toc}{fg=qaaccent}
\setbeamercolor{subsection in toc}{fg=qamuted}
% Sommaire de debut de section : section en cours en violet et gras, les autres en gris
\setbeamercolor{section in toc shaded}{fg=qamuted}
\setbeamercolor{subsection in toc shaded}{fg=qamuted}

\setbeamertemplate{blocks}[rounded][shadow=false]
\setbeamertemplate{itemize item}{\raisebox{0.15em}{\scriptsize\textcolor{qaaccent}{\ensuremath{\blacktriangleright}}}}
\setbeamertemplate{itemize subitem}{\raisebox{0.1em}{\tiny\textcolor{qaaccent2}{\ensuremath{\bullet}}}}

\setbeamerfont{frametitle}{series=\bfseries, size=\Large}
\setbeamerfont{title}{series=\bfseries, size=\Large}
\setbeamerfont{subtitle}{size=\normalsize}
\setbeamerfont{institute}{size=\scriptsize}
\setbeamerfont{date}{size=\scriptsize}
\setbeamerfont{author}{size=\small}

\metroset{sectionpage=none, subsectionpage=none, numbering=fraction, progressbar=frametitle}
\setbeamertemplate{section in toc}[sections numbered]

% ---------------------------------------------------------------
% LISTINGS -- theme sombre (lstlisting reste utilisable tel quel)
% \lstcode : police + couleur de base, a reutiliser dans les
%            surcharges locales (ex. basicstyle=\lstcode\scriptsize)
% ---------------------------------------------------------------
\newcommand{\lstcode}{\ttfamily\color{qatext}}

\lstdefinelanguage{TypeScript}{
  keywords={typeof, new, true, false, catch, function, return, null, switch, var,
    if, in, while, do, else, case, break, const, let, interface, class, extends,
    implements, enum, export, import, this, super, constructor, public, private,
    protected, static, async, await, throw, throws, try, finally, void},
  ndkeywords={boolean, any, number, string, int, long, double},
  sensitive=true,
  comment=[l]{//},
  morecomment=[s]{/*}{*/},
  morestring=[b]',
  morestring=[b]"
}

\lstset{
  language=Java,
  basicstyle=\lstcode\footnotesize,
  keywordstyle=\color{qacodekw}\bfseries,
  ndkeywordstyle=\color{qaaccent2}\bfseries,
  identifierstyle=\color{qatext},
  stringstyle=\color{qacodestr},
  commentstyle=\color{qacodecom}\itshape,
  morecomment=[s][\color{qacodeann}]{/**}{*/},
  backgroundcolor=\color{qapanelalt},
  rulecolor=\color{qaline},
  frame=single,
  framesep=3pt,
  xleftmargin=4pt,
  xrightmargin=4pt,
  keepspaces=true,
  columns=fullflexible,
  showstringspaces=false,
  breaklines=true,
  tabsize=2,
  escapeinside={(*@}{@*)},
  % accents UTF-8 dans le code (listings ne gere pas l'UTF-8 multi-octets)
  literate={é}{{\'e}}1 {è}{{\`e}}1 {ê}{{\^e}}1 {à}{{\`a}}1 {â}{{\^a}}1
           {ù}{{\`u}}1 {û}{{\^u}}1 {ô}{{\^o}}1 {î}{{\^i}}1 {ç}{{\c c}}1
           {É}{{\'E}}1 {À}{{\`A}}1 {€}{{\euro{}}}1 {…}{{\ldots}}1,
}

% ---------------------------------------------------------------
% CARD-STYLE BOXES
% ---------------------------------------------------------------
\tcbset{
  figcard/.style={colback=white, colframe=white, boxrule=0pt, arc=4pt,
    left=3pt, right=3pt, top=3pt, bottom=3pt, boxsep=0pt,
    halign=center, valign=center},
  datacard/.style={colback=qapanel, colframe=qaline, boxrule=0.4pt, arc=3pt,
    left=8pt, right=8pt, top=6pt, bottom=6pt, coltext=qatext,
    before upper={\color{qatext}}, halign=center},
  codebox/.style={colback=qapanelalt, colframe=qaline, boxrule=0.4pt, arc=3pt,
    left=8pt, right=8pt, top=5pt, bottom=5pt, listing only,
    listing options={basicstyle=\ttfamily\small\color{qaaccent},
      breaklines=true, showstringspaces=false, columns=fullflexible,
      keepspaces=true, frame=none, backgroundcolor={}}}
}
\newtcolorbox{figcard}[1][]{figcard,#1}
\newtcolorbox{datacard}{datacard}
\newtcblisting{codebox}{codebox}

% ---------------------------------------------------------------
% HANDY MACROS
% ---------------------------------------------------------------
% Mention d'elaboration (discrete : plus sombre que le texte secondaire)
\definecolor{qaaimention}{HTML}{4E5566}
\newcommand{\raimention}{{\tiny\color{qaaimention}Supports \'elabor\'es avec l'aide de Claude Code Opus 5.5}}
\newcommand{\hl}[1]{\alert{#1}}
\newcommand{\source}[1]{{\scriptsize\color{qamuted}#1}}

% Sommaire automatique en debut de section
\AtBeginSection[]{%
  \begin{frame}<beamer>{Plan du cours}
    \tableofcontents[currentsection]
  \end{frame}}

% Page de titre : \rtitleslide{numero du cours}{titre du cours}
% Les auteurs sont repris de \rauthors (modifiable avant \begin{document})
\newcommand{\rauthors}{Chouki Tibermacine}
\newcommand{\rtitleslide}[2]{%
\begin{frame}[plain,noframenumbering]
  \begin{center}
    \vspace{2mm}
    {\color{qaaccent}\rule{2.1cm}{1.2pt}}\\[6mm]
    {\Large\bfseries\color{qatext} R4.02}\\[1mm]
    {\Large\bfseries\color{qatext} Qualit\'e de d\'eveloppement}\\[5mm]
    {\large\color{qamuted} Cours #1}\\[1mm]
    {\large\color{qaaccent} #2}\\[6mm]
    \includegraphics[height=2.1cm]{iut_logo_transparent.png}

    \vspace{4mm}

    {\small\color{qatext} \rauthors}\\[1mm]
    {\scriptsize\color{qamuted} BUT 2 Informatique --- IUT de Vannes --- Universit\'e Bretagne Sud}\\[2mm]
    \raimention
  \end{center}
\end{frame}
}

% ---- Diapo de demonstration : \rdemo{numero}{titre}{contenu (items)}{fichier(s) a ouvrir} ----
\newcommand{\rdemo}[4]{%
\begin{frame}{D\'emo #1 --- #2}
  \begin{tikzpicture}[remember picture, overlay]
    \node[anchor=north east, fill=qaaccent2, text=qabg, font=\scriptsize\bfseries,
      rounded corners=3pt, inner sep=4pt] at ([xshift=-0.4cm, yshift=-0.35cm]current page.north east) {D\'EMO};
  \end{tikzpicture}
  \begin{columns}[T]
    \begin{column}{0.62\textwidth}
      \begin{itemize}
        #3
      \end{itemize}
    \end{column}
    \begin{column}{0.34\textwidth}
      \begin{datacard}
        \raggedright\scriptsize
        {\color{qamuted} Dans l'IDE}\\[1mm]
        #4
      \end{datacard}
    \end{column}
  \end{columns}
\end{frame}}

% ---- Diapo de fin « Questions ? » ----
\newcommand{\rcontact}{chouki.tibermacine@univ-ubs.fr}
\newcommand{\rquestionslide}{%
\begin{frame}{Questions ?}
  \begin{center}
    \vspace{2mm}
    \begin{tikzpicture}
      \node[regular polygon, regular polygon sides=6, draw=qaline, thick,
        minimum size=4.6cm, rotate=90] (hexin) {};
      \node[regular polygon, regular polygon sides=6, draw=qaaccent, dashed,
        thick, minimum size=6.2cm, rotate=90] (hexout) {};
      \node[font=\fontsize{92}{92}\selectfont\bfseries, color=qaaccent2] at (hexin) {?};
    \end{tikzpicture}\\[5mm]
    {\small\color{qamuted} \rcontact}
  \end{center}
\end{frame}}

%   \graphicspath{{../../latex/common/assets/}{img/}{figs/}}
% ---------------------------------------------------------------
\usetheme{metropolis}
\usepackage{appendixnumberbeamer}

\usepackage[T1]{fontenc}
\usepackage[utf8]{inputenc}
\usepackage[french]{babel}
\usepackage[scale=0.94]{tgheros}
\renewcommand{\familydefault}{\sfdefault}
\usepackage{amsmath,amssymb,amsfonts}
\usepackage{graphicx}
\usepackage{booktabs}
\usepackage{multirow}
\usepackage{tikz}
\usetikzlibrary{shapes.geometric, shapes, arrows, arrows.meta, positioning, calc, fit}
\usepackage{xcolor}
\usepackage{listings}
\usepackage[skins,listings]{tcolorbox}
\usepackage{qrcode}
\usepackage{eurosym}
\usepackage{hyperref}

% ---------------------------------------------------------------
% PALETTE (dark-mode tech: near-black canvas, violet (logo IUT) + gold accents)
% ---------------------------------------------------------------
\definecolor{qabg}{HTML}{0E1016}      % canvas background
\definecolor{qapanel}{HTML}{181B24}   % card / block panel
\definecolor{qapanelalt}{HTML}{20242F}% secondary panel tint
\definecolor{qaline}{HTML}{2B303D}    % hairline rules / table rules
\definecolor{qatext}{HTML}{ECEEF2}    % primary text (off-white)
\definecolor{qamuted}{HTML}{8B93A6}   % secondary / caption text
\definecolor{qaaccent}{HTML}{D93ED7}  % primary accent -- violet du logo IUT, eclairci pour fond sombre
\definecolor{qaaccent2}{HTML}{F2B84B} % secondary accent -- or/ambre
\definecolor{qacodekw}{HTML}{C792EA}  % code : mots-cles
\definecolor{qacodestr}{HTML}{C3E88D} % code : chaines
\definecolor{qacodecom}{HTML}{7F8AA3} % code : commentaires
\definecolor{qacodeann}{HTML}{F2B84B} % code : annotations / javadoc

\setbeamercolor{background canvas}{bg=qabg}
\setbeamercolor{normal text}{fg=qatext}
\setbeamercolor{palette primary}{bg=qabg, fg=qamuted}
\setbeamercolor{frametitle}{bg=qabg, fg=qatext}
\setbeamercolor{progress bar}{fg=qaaccent, bg=qaline}
\setbeamercolor{alerted text}{fg=qaaccent2}
\setbeamercolor{block title}{bg=qapanel, fg=qaaccent}
\setbeamercolor{block body}{bg=qapanel, fg=qatext}
\setbeamercolor{block title alerted}{bg=qapanel, fg=qaaccent2}
\setbeamercolor{block body alerted}{bg=qapanel, fg=qatext}
\setbeamercolor{block title example}{bg=qapanelalt, fg=qaaccent}
\setbeamercolor{block body example}{bg=qapanelalt, fg=qatext}
\setbeamercolor{example text}{fg=qaaccent}
\setbeamercolor{itemize item}{fg=qaaccent}
\setbeamercolor{itemize subitem}{fg=qaaccent2}
\setbeamercolor{enumerate item}{fg=qaaccent}
\setbeamercolor{enumerate subitem}{fg=qaaccent2}
\setbeamercolor{caption}{fg=qamuted}
\setbeamercolor{page number in head/foot}{fg=qamuted, bg=qabg}
\setbeamercolor{title}{fg=qatext}
\setbeamercolor{subtitle}{fg=qamuted}
\setbeamercolor{author}{fg=qatext}
\setbeamercolor{institute}{fg=qamuted}
\setbeamercolor{date}{fg=qamuted}
\setbeamercolor{section in toc}{fg=qaaccent}
\setbeamercolor{subsection in toc}{fg=qamuted}
% Sommaire de debut de section : section en cours en violet et gras, les autres en gris
\setbeamercolor{section in toc shaded}{fg=qamuted}
\setbeamercolor{subsection in toc shaded}{fg=qamuted}

\setbeamertemplate{blocks}[rounded][shadow=false]
\setbeamertemplate{itemize item}{\raisebox{0.15em}{\scriptsize\textcolor{qaaccent}{\ensuremath{\blacktriangleright}}}}
\setbeamertemplate{itemize subitem}{\raisebox{0.1em}{\tiny\textcolor{qaaccent2}{\ensuremath{\bullet}}}}

\setbeamerfont{frametitle}{series=\bfseries, size=\Large}
\setbeamerfont{title}{series=\bfseries, size=\Large}
\setbeamerfont{subtitle}{size=\normalsize}
\setbeamerfont{institute}{size=\scriptsize}
\setbeamerfont{date}{size=\scriptsize}
\setbeamerfont{author}{size=\small}

\metroset{sectionpage=none, subsectionpage=none, numbering=fraction, progressbar=frametitle}
\setbeamertemplate{section in toc}[sections numbered]

% ---------------------------------------------------------------
% LISTINGS -- theme sombre (lstlisting reste utilisable tel quel)
% \lstcode : police + couleur de base, a reutiliser dans les
%            surcharges locales (ex. basicstyle=\lstcode\scriptsize)
% ---------------------------------------------------------------
\newcommand{\lstcode}{\ttfamily\color{qatext}}

\lstdefinelanguage{TypeScript}{
  keywords={typeof, new, true, false, catch, function, return, null, switch, var,
    if, in, while, do, else, case, break, const, let, interface, class, extends,
    implements, enum, export, import, this, super, constructor, public, private,
    protected, static, async, await, throw, throws, try, finally, void},
  ndkeywords={boolean, any, number, string, int, long, double},
  sensitive=true,
  comment=[l]{//},
  morecomment=[s]{/*}{*/},
  morestring=[b]',
  morestring=[b]"
}

\lstset{
  language=Java,
  basicstyle=\lstcode\footnotesize,
  keywordstyle=\color{qacodekw}\bfseries,
  ndkeywordstyle=\color{qaaccent2}\bfseries,
  identifierstyle=\color{qatext},
  stringstyle=\color{qacodestr},
  commentstyle=\color{qacodecom}\itshape,
  morecomment=[s][\color{qacodeann}]{/**}{*/},
  backgroundcolor=\color{qapanelalt},
  rulecolor=\color{qaline},
  frame=single,
  framesep=3pt,
  xleftmargin=4pt,
  xrightmargin=4pt,
  keepspaces=true,
  columns=fullflexible,
  showstringspaces=false,
  breaklines=true,
  tabsize=2,
  escapeinside={(*@}{@*)},
  % accents UTF-8 dans le code (listings ne gere pas l'UTF-8 multi-octets)
  literate={é}{{\'e}}1 {è}{{\`e}}1 {ê}{{\^e}}1 {à}{{\`a}}1 {â}{{\^a}}1
           {ù}{{\`u}}1 {û}{{\^u}}1 {ô}{{\^o}}1 {î}{{\^i}}1 {ç}{{\c c}}1
           {É}{{\'E}}1 {À}{{\`A}}1 {€}{{\euro{}}}1 {…}{{\ldots}}1,
}

% ---------------------------------------------------------------
% CARD-STYLE BOXES
% ---------------------------------------------------------------
\tcbset{
  figcard/.style={colback=white, colframe=white, boxrule=0pt, arc=4pt,
    left=3pt, right=3pt, top=3pt, bottom=3pt, boxsep=0pt,
    halign=center, valign=center},
  datacard/.style={colback=qapanel, colframe=qaline, boxrule=0.4pt, arc=3pt,
    left=8pt, right=8pt, top=6pt, bottom=6pt, coltext=qatext,
    before upper={\color{qatext}}, halign=center},
  codebox/.style={colback=qapanelalt, colframe=qaline, boxrule=0.4pt, arc=3pt,
    left=8pt, right=8pt, top=5pt, bottom=5pt, listing only,
    listing options={basicstyle=\ttfamily\small\color{qaaccent},
      breaklines=true, showstringspaces=false, columns=fullflexible,
      keepspaces=true, frame=none, backgroundcolor={}}}
}
\newtcolorbox{figcard}[1][]{figcard,#1}
\newtcolorbox{datacard}{datacard}
\newtcblisting{codebox}{codebox}

% ---------------------------------------------------------------
% HANDY MACROS
% ---------------------------------------------------------------
% Mention d'elaboration (discrete : plus sombre que le texte secondaire)
\definecolor{qaaimention}{HTML}{4E5566}
\newcommand{\raimention}{{\tiny\color{qaaimention}Supports \'elabor\'es avec l'aide de Claude Code Opus 5.5}}
\newcommand{\hl}[1]{\alert{#1}}
\newcommand{\source}[1]{{\scriptsize\color{qamuted}#1}}

% Sommaire automatique en debut de section
\AtBeginSection[]{%
  \begin{frame}<beamer>{Plan du cours}
    \tableofcontents[currentsection]
  \end{frame}}

% Page de titre : \rtitleslide{numero du cours}{titre du cours}
% Les auteurs sont repris de \rauthors (modifiable avant \begin{document})
\newcommand{\rauthors}{Chouki Tibermacine}
\newcommand{\rtitleslide}[2]{%
\begin{frame}[plain,noframenumbering]
  \begin{center}
    \vspace{2mm}
    {\color{qaaccent}\rule{2.1cm}{1.2pt}}\\[6mm]
    {\Large\bfseries\color{qatext} R4.02}\\[1mm]
    {\Large\bfseries\color{qatext} Qualit\'e de d\'eveloppement}\\[5mm]
    {\large\color{qamuted} Cours #1}\\[1mm]
    {\large\color{qaaccent} #2}\\[6mm]
    \includegraphics[height=2.1cm]{iut_logo_transparent.png}

    \vspace{4mm}

    {\small\color{qatext} \rauthors}\\[1mm]
    {\scriptsize\color{qamuted} BUT 2 Informatique --- IUT de Vannes --- Universit\'e Bretagne Sud}\\[2mm]
    \raimention
  \end{center}
\end{frame}
}

% ---- Diapo de demonstration : \rdemo{numero}{titre}{contenu (items)}{fichier(s) a ouvrir} ----
\newcommand{\rdemo}[4]{%
\begin{frame}{D\'emo #1 --- #2}
  \begin{tikzpicture}[remember picture, overlay]
    \node[anchor=north east, fill=qaaccent2, text=qabg, font=\scriptsize\bfseries,
      rounded corners=3pt, inner sep=4pt] at ([xshift=-0.4cm, yshift=-0.35cm]current page.north east) {D\'EMO};
  \end{tikzpicture}
  \begin{columns}[T]
    \begin{column}{0.62\textwidth}
      \begin{itemize}
        #3
      \end{itemize}
    \end{column}
    \begin{column}{0.34\textwidth}
      \begin{datacard}
        \raggedright\scriptsize
        {\color{qamuted} Dans l'IDE}\\[1mm]
        #4
      \end{datacard}
    \end{column}
  \end{columns}
\end{frame}}

% ---- Diapo de fin « Questions ? » ----
\newcommand{\rcontact}{chouki.tibermacine@univ-ubs.fr}
\newcommand{\rquestionslide}{%
\begin{frame}{Questions ?}
  \begin{center}
    \vspace{2mm}
    \begin{tikzpicture}
      \node[regular polygon, regular polygon sides=6, draw=qaline, thick,
        minimum size=4.6cm, rotate=90] (hexin) {};
      \node[regular polygon, regular polygon sides=6, draw=qaaccent, dashed,
        thick, minimum size=6.2cm, rotate=90] (hexout) {};
      \node[font=\fontsize{92}{92}\selectfont\bfseries, color=qaaccent2] at (hexin) {?};
    \end{tikzpicture}\\[5mm]
    {\small\color{qamuted} \rcontact}
  \end{center}
\end{frame}}

%   \graphicspath{{../../latex/common/assets/}{img/}{figs/}}
% ---------------------------------------------------------------
\usetheme{metropolis}
\usepackage{appendixnumberbeamer}

\usepackage[T1]{fontenc}
\usepackage[utf8]{inputenc}
\usepackage[french]{babel}
\usepackage[scale=0.94]{tgheros}
\renewcommand{\familydefault}{\sfdefault}
\usepackage{amsmath,amssymb,amsfonts}
\usepackage{graphicx}
\usepackage{booktabs}
\usepackage{multirow}
\usepackage{tikz}
\usetikzlibrary{shapes.geometric, shapes, arrows, arrows.meta, positioning, calc, fit}
\usepackage{xcolor}
\usepackage{listings}
\usepackage[skins,listings]{tcolorbox}
\usepackage{qrcode}
\usepackage{eurosym}
\usepackage{hyperref}

% ---------------------------------------------------------------
% PALETTE (dark-mode tech: near-black canvas, violet (logo IUT) + gold accents)
% ---------------------------------------------------------------
\definecolor{qabg}{HTML}{0E1016}      % canvas background
\definecolor{qapanel}{HTML}{181B24}   % card / block panel
\definecolor{qapanelalt}{HTML}{20242F}% secondary panel tint
\definecolor{qaline}{HTML}{2B303D}    % hairline rules / table rules
\definecolor{qatext}{HTML}{ECEEF2}    % primary text (off-white)
\definecolor{qamuted}{HTML}{8B93A6}   % secondary / caption text
\definecolor{qaaccent}{HTML}{D93ED7}  % primary accent -- violet du logo IUT, eclairci pour fond sombre
\definecolor{qaaccent2}{HTML}{F2B84B} % secondary accent -- or/ambre
\definecolor{qacodekw}{HTML}{C792EA}  % code : mots-cles
\definecolor{qacodestr}{HTML}{C3E88D} % code : chaines
\definecolor{qacodecom}{HTML}{7F8AA3} % code : commentaires
\definecolor{qacodeann}{HTML}{F2B84B} % code : annotations / javadoc

\setbeamercolor{background canvas}{bg=qabg}
\setbeamercolor{normal text}{fg=qatext}
\setbeamercolor{palette primary}{bg=qabg, fg=qamuted}
\setbeamercolor{frametitle}{bg=qabg, fg=qatext}
\setbeamercolor{progress bar}{fg=qaaccent, bg=qaline}
\setbeamercolor{alerted text}{fg=qaaccent2}
\setbeamercolor{block title}{bg=qapanel, fg=qaaccent}
\setbeamercolor{block body}{bg=qapanel, fg=qatext}
\setbeamercolor{block title alerted}{bg=qapanel, fg=qaaccent2}
\setbeamercolor{block body alerted}{bg=qapanel, fg=qatext}
\setbeamercolor{block title example}{bg=qapanelalt, fg=qaaccent}
\setbeamercolor{block body example}{bg=qapanelalt, fg=qatext}
\setbeamercolor{example text}{fg=qaaccent}
\setbeamercolor{itemize item}{fg=qaaccent}
\setbeamercolor{itemize subitem}{fg=qaaccent2}
\setbeamercolor{enumerate item}{fg=qaaccent}
\setbeamercolor{enumerate subitem}{fg=qaaccent2}
\setbeamercolor{caption}{fg=qamuted}
\setbeamercolor{page number in head/foot}{fg=qamuted, bg=qabg}
\setbeamercolor{title}{fg=qatext}
\setbeamercolor{subtitle}{fg=qamuted}
\setbeamercolor{author}{fg=qatext}
\setbeamercolor{institute}{fg=qamuted}
\setbeamercolor{date}{fg=qamuted}
\setbeamercolor{section in toc}{fg=qaaccent}
\setbeamercolor{subsection in toc}{fg=qamuted}
% Sommaire de debut de section : section en cours en violet et gras, les autres en gris
\setbeamercolor{section in toc shaded}{fg=qamuted}
\setbeamercolor{subsection in toc shaded}{fg=qamuted}

\setbeamertemplate{blocks}[rounded][shadow=false]
\setbeamertemplate{itemize item}{\raisebox{0.15em}{\scriptsize\textcolor{qaaccent}{\ensuremath{\blacktriangleright}}}}
\setbeamertemplate{itemize subitem}{\raisebox{0.1em}{\tiny\textcolor{qaaccent2}{\ensuremath{\bullet}}}}

\setbeamerfont{frametitle}{series=\bfseries, size=\Large}
\setbeamerfont{title}{series=\bfseries, size=\Large}
\setbeamerfont{subtitle}{size=\normalsize}
\setbeamerfont{institute}{size=\scriptsize}
\setbeamerfont{date}{size=\scriptsize}
\setbeamerfont{author}{size=\small}

\metroset{sectionpage=none, subsectionpage=none, numbering=fraction, progressbar=frametitle}
\setbeamertemplate{section in toc}[sections numbered]

% ---------------------------------------------------------------
% LISTINGS -- theme sombre (lstlisting reste utilisable tel quel)
% \lstcode : police + couleur de base, a reutiliser dans les
%            surcharges locales (ex. basicstyle=\lstcode\scriptsize)
% ---------------------------------------------------------------
\newcommand{\lstcode}{\ttfamily\color{qatext}}

\lstdefinelanguage{TypeScript}{
  keywords={typeof, new, true, false, catch, function, return, null, switch, var,
    if, in, while, do, else, case, break, const, let, interface, class, extends,
    implements, enum, export, import, this, super, constructor, public, private,
    protected, static, async, await, throw, throws, try, finally, void},
  ndkeywords={boolean, any, number, string, int, long, double},
  sensitive=true,
  comment=[l]{//},
  morecomment=[s]{/*}{*/},
  morestring=[b]',
  morestring=[b]"
}

\lstset{
  language=Java,
  basicstyle=\lstcode\footnotesize,
  keywordstyle=\color{qacodekw}\bfseries,
  ndkeywordstyle=\color{qaaccent2}\bfseries,
  identifierstyle=\color{qatext},
  stringstyle=\color{qacodestr},
  commentstyle=\color{qacodecom}\itshape,
  morecomment=[s][\color{qacodeann}]{/**}{*/},
  backgroundcolor=\color{qapanelalt},
  rulecolor=\color{qaline},
  frame=single,
  framesep=3pt,
  xleftmargin=4pt,
  xrightmargin=4pt,
  keepspaces=true,
  columns=fullflexible,
  showstringspaces=false,
  breaklines=true,
  tabsize=2,
  escapeinside={(*@}{@*)},
  % accents UTF-8 dans le code (listings ne gere pas l'UTF-8 multi-octets)
  literate={é}{{\'e}}1 {è}{{\`e}}1 {ê}{{\^e}}1 {à}{{\`a}}1 {â}{{\^a}}1
           {ù}{{\`u}}1 {û}{{\^u}}1 {ô}{{\^o}}1 {î}{{\^i}}1 {ç}{{\c c}}1
           {É}{{\'E}}1 {À}{{\`A}}1 {€}{{\euro{}}}1 {…}{{\ldots}}1,
}

% ---------------------------------------------------------------
% CARD-STYLE BOXES
% ---------------------------------------------------------------
\tcbset{
  figcard/.style={colback=white, colframe=white, boxrule=0pt, arc=4pt,
    left=3pt, right=3pt, top=3pt, bottom=3pt, boxsep=0pt,
    halign=center, valign=center},
  datacard/.style={colback=qapanel, colframe=qaline, boxrule=0.4pt, arc=3pt,
    left=8pt, right=8pt, top=6pt, bottom=6pt, coltext=qatext,
    before upper={\color{qatext}}, halign=center},
  codebox/.style={colback=qapanelalt, colframe=qaline, boxrule=0.4pt, arc=3pt,
    left=8pt, right=8pt, top=5pt, bottom=5pt, listing only,
    listing options={basicstyle=\ttfamily\small\color{qaaccent},
      breaklines=true, showstringspaces=false, columns=fullflexible,
      keepspaces=true, frame=none, backgroundcolor={}}}
}
\newtcolorbox{figcard}[1][]{figcard,#1}
\newtcolorbox{datacard}{datacard}
\newtcblisting{codebox}{codebox}

% ---------------------------------------------------------------
% HANDY MACROS
% ---------------------------------------------------------------
% Mention d'elaboration (discrete : plus sombre que le texte secondaire)
\definecolor{qaaimention}{HTML}{4E5566}
\newcommand{\raimention}{{\tiny\color{qaaimention}Supports \'elabor\'es avec l'aide de Claude Code Opus 5.5}}
\newcommand{\hl}[1]{\alert{#1}}
\newcommand{\source}[1]{{\scriptsize\color{qamuted}#1}}

% Sommaire automatique en debut de section
\AtBeginSection[]{%
  \begin{frame}<beamer>{Plan du cours}
    \tableofcontents[currentsection]
  \end{frame}}

% Page de titre : \rtitleslide{numero du cours}{titre du cours}
% Les auteurs sont repris de \rauthors (modifiable avant \begin{document})
\newcommand{\rauthors}{Chouki Tibermacine}
\newcommand{\rtitleslide}[2]{%
\begin{frame}[plain,noframenumbering]
  \begin{center}
    \vspace{2mm}
    {\color{qaaccent}\rule{2.1cm}{1.2pt}}\\[6mm]
    {\Large\bfseries\color{qatext} R4.02}\\[1mm]
    {\Large\bfseries\color{qatext} Qualit\'e de d\'eveloppement}\\[5mm]
    {\large\color{qamuted} Cours #1}\\[1mm]
    {\large\color{qaaccent} #2}\\[6mm]
    \includegraphics[height=2.1cm]{iut_logo_transparent.png}

    \vspace{4mm}

    {\small\color{qatext} \rauthors}\\[1mm]
    {\scriptsize\color{qamuted} BUT 2 Informatique --- IUT de Vannes --- Universit\'e Bretagne Sud}\\[2mm]
    \raimention
  \end{center}
\end{frame}
}

% ---- Diapo de demonstration : \rdemo{numero}{titre}{contenu (items)}{fichier(s) a ouvrir} ----
\newcommand{\rdemo}[4]{%
\begin{frame}{D\'emo #1 --- #2}
  \begin{tikzpicture}[remember picture, overlay]
    \node[anchor=north east, fill=qaaccent2, text=qabg, font=\scriptsize\bfseries,
      rounded corners=3pt, inner sep=4pt] at ([xshift=-0.4cm, yshift=-0.35cm]current page.north east) {D\'EMO};
  \end{tikzpicture}
  \begin{columns}[T]
    \begin{column}{0.62\textwidth}
      \begin{itemize}
        #3
      \end{itemize}
    \end{column}
    \begin{column}{0.34\textwidth}
      \begin{datacard}
        \raggedright\scriptsize
        {\color{qamuted} Dans l'IDE}\\[1mm]
        #4
      \end{datacard}
    \end{column}
  \end{columns}
\end{frame}}

% ---- Diapo de fin « Questions ? » ----
\newcommand{\rcontact}{chouki.tibermacine@univ-ubs.fr}
\newcommand{\rquestionslide}{%
\begin{frame}{Questions ?}
  \begin{center}
    \vspace{2mm}
    \begin{tikzpicture}
      \node[regular polygon, regular polygon sides=6, draw=qaline, thick,
        minimum size=4.6cm, rotate=90] (hexin) {};
      \node[regular polygon, regular polygon sides=6, draw=qaaccent, dashed,
        thick, minimum size=6.2cm, rotate=90] (hexout) {};
      \node[font=\fontsize{92}{92}\selectfont\bfseries, color=qaaccent2] at (hexin) {?};
    \end{tikzpicture}\\[5mm]
    {\small\color{qamuted} \rcontact}
  \end{center}
\end{frame}}

\graphicspath{{../../latex/common/assets/}{img/}}

\title{Tests d'intégration avec une base de données}
\author[C. Tibermacine]{Chouki Tibermacine}
\date{}

\lstdefinestyle{sortie}{language={}, basicstyle=\lstcode\scriptsize, keywordstyle={}, commentstyle={}, stringstyle={}}

\tikzset{
  comp/.style={rectangle, rounded corners=2pt, draw=qaaccent, thick, fill=qapanel, text=qatext,
    minimum height=0.55cm, minimum width=2.5cm, align=center, font=\scriptsize},
  dep/.style={rectangle, rounded corners=2pt, draw=qamuted, thick, fill=qapanelalt, text=qatext,
    minimum height=0.55cm, minimum width=2.5cm, align=center, font=\scriptsize},
  dbl/.style={rectangle, rounded corners=2pt, draw=qaaccent2, thick, dashed, fill=qapanelalt, text=qatext,
    minimum height=0.55cm, minimum width=2.5cm, align=center, font=\scriptsize},
  arr/.style={-Latex, qamuted, thick}
}

\begin{document}

\rtitleslide{4}{Tests d'intégration avec une base de données}

\begin{frame}{Où en est-on ?}
  \begin{columns}[c]
    \begin{column}{0.48\textwidth}
      \begin{itemize}
        \item CM1 et CM3 : des tests \textbf{unitaires}, rapides, isolés par des doublures
        \item CM2 : des tests d'\textbf{acceptation}, écrits avec le client
        \item Mais avec des doublures, on a seulement \textbf{supposé} que la base, le SQL et la configuration se comportent comme prévu
        \item Aujourd'hui : vérifier que les morceaux \textbf{s'assemblent} vraiment
      \end{itemize}
    \end{column}
    \begin{column}{0.48\textwidth}
      \begin{center}
      \begin{tikzpicture}[x=0.8cm, y=1cm]
        \fill[qaaccent!35!qapanel] (0,0) -- (6,0) -- (4.6,1.4) -- (1.4,1.4) -- cycle;
        \fill[qaaccent2] (1.4,1.4) -- (4.6,1.4) -- (3.7,2.3) -- (2.3,2.3) -- cycle;
        \fill[qaaccent!35!qapanel] (2.3,2.3) -- (3.7,2.3) -- (3,3.0) -- cycle;
        \node[font=\small\bfseries, text=qatext] at (3,0.7) {Unitaires};
        \node[font=\scriptsize\bfseries, text=qabg] at (3,1.85) {Intégration};
        \node[font=\tiny\bfseries, text=qatext] at (3,2.55) {E2E};
        \draw[arr, qaaccent2] (6.6,1.85) -- (4.4,1.85);
        \node[font=\scriptsize, text=qaaccent2, anchor=west] at (6.6,1.85) {CM4};
      \end{tikzpicture}
      \end{center}
    \end{column}
  \end{columns}
\end{frame}

% ================================================================
\section{Qu'est-ce qu'on intègre ?}

\begin{frame}{Ce que les tests unitaires ne voient pas}
  \begin{columns}[T]
    \begin{column}{0.48\textwidth}
      \begin{block}{Côté base de données}
        \begin{itemize}
          \item une requête SQL ou JPQL mal écrite
          \item un mapping JPA faux (nom de colonne, type, relation)
          \item une contrainte (\texttt{UNIQUE}, \texttt{NOT NULL}, clé étrangère) qui rejette les données
          \item une transaction qui laisse des données incohérentes
          \item un trigger ou une procédure stockée (R3.07)
        \end{itemize}
      \end{block}
    \end{column}
    \begin{column}{0.48\textwidth}
      \begin{block}{Côté application}
        \begin{itemize}
          \item une configuration Spring incomplète (un bean manquant)
          \item une URL ou un format JSON mal défini dans un contrôleur
          \item une sérialisation de \texttt{BigDecimal} ou de date inattendue
          \item un mock qui renvoie ce que la vraie base ne renverrait jamais
        \end{itemize}
      \end{block}
    \end{column}
  \end{columns}
  \vspace{3mm}
  \begin{center}
    Tous ces bugs passent les tests unitaires\ldots{} et apparaissent au premier démarrage réel.
  \end{center}
\end{frame}

\begin{frame}{Qu'est-ce qu'un test d'intégration ?}
  \begin{columns}[c]
    \begin{column}{0.5\textwidth}
      \begin{center}
      \begin{tikzpicture}[node distance=0.35cm]
        \node[comp] (ctrl) {Contrôleur REST};
        \node[comp, below=of ctrl] (svc) {Service};
        \node[comp, below=of svc] (repo) {Repository};
        \node[comp, below=of repo, draw=qaaccent2] (db) {PostgreSQL};
        \node[dbl, right=0.7cm of svc] (api) {API de change\\{\tiny doublure}};
        \draw[arr] (ctrl) -- (svc);
        \draw[arr] (svc) -- (repo);
        \draw[arr] (repo) -- (db);
        \draw[arr] (svc) -- (api);
        \draw[qaaccent2, thick, rounded corners, dashed] ([xshift=-0.25cm, yshift=0.2cm]ctrl.north west) rectangle ([xshift=0.25cm, yshift=-0.2cm]db.south east);
        \node[font=\tiny, text=qaaccent2, anchor=north] at ([yshift=-0.25cm]db.south) {composants réels, testés ensemble};
      \end{tikzpicture}
      \end{center}
    \end{column}
    \begin{column}{0.48\textwidth}
      \begin{itemize}
        \item Plusieurs composants \textbf{réels} testés \textbf{ensemble}, à travers leurs interfaces
        \item Les plus fréquents :
          \begin{itemize}
            \item code + base de données
            \item code + API externe (adaptateur HTTP)
            \item toutes les couches de l'application
          \end{itemize}
        \item On garde des doublures pour ce qui est \textbf{hors} du périmètre testé (ici : l'API de change)
        \item Plus lents que les tests unitaires (secondes, pas millisecondes), donc moins nombreux
      \end{itemize}
    \end{column}
  \end{columns}
\end{frame}

\begin{frame}{Stratégies d'intégration}
  \begin{columns}[T]
    \begin{column}{0.32\textwidth}
      \begin{block}{Big bang}
        \small
        Tout assembler d'un coup, puis tester.
        \begin{itemize}
          \item simple à organiser
          \item \alert{si ça échoue, où est le problème ?}
        \end{itemize}
      \end{block}
    \end{column}
    \begin{column}{0.32\textwidth}
      \begin{block}{Ascendante}
        \small
        Des couches basses vers les hautes : repository + base, puis service, puis contrôleur.
        \begin{itemize}
          \item chaque étape repose sur une couche déjà vérifiée
          \item l'interface utilisateur arrive en dernier
        \end{itemize}
      \end{block}
    \end{column}
    \begin{column}{0.32\textwidth}
      \begin{block}{Descendante}
        \small
        Des couches hautes vers les basses, en remplaçant le bas par des doublures.
        \begin{itemize}
          \item on montre tôt un parcours complet
          \item beaucoup de doublures à écrire
        \end{itemize}
      \end{block}
    \end{column}
  \end{columns}
  \vspace{3mm}
  \begin{itemize}
    \item En pratique, avec une intégration continue, on intègre \textbf{petit à petit} et on teste à chaque étape : ni big bang, ni grand plan
    \item Dans ce cours : ascendant (repository + base, puis toute l'application)
  \end{itemize}
\end{frame}

% ================================================================
\section{Quelle base de données pour les tests ?}

\begin{frame}{Pourquoi pas la base de développement ?}
  \begin{itemize}
    \item Les tests \textbf{écrivent} : ils créent des comptes, font des virements, suppriment\ldots
    \item Avec une base partagée :
      \begin{itemize}
        \item les tests dépendent des données déjà présentes (« ça passe chez moi »)
        \item ils dépendent de l'\textbf{ordre} d'exécution : un test laisse des données qui en font échouer un autre
        \item deux développeurs qui lancent les tests en même temps se gênent
        \item vos données de développement sont modifiées ou effacées
      \end{itemize}
    \item Il faut une base \textbf{dédiée aux tests}, dans un état \textbf{connu} au début de chaque test
    \item Et en CI (CM6), il n'y a pas de base de développement du tout
  \end{itemize}
  \vspace{2mm}
  \begin{center}
    Deux options : une base \textbf{en mémoire} (H2) ou une \textbf{vraie} base dans un conteneur (Testcontainers)
  \end{center}
\end{frame}

\begin{frame}{Option 1 : H2, une base en mémoire}
  \small
  \begin{columns}[T]
    \begin{column}{0.48\textwidth}
      \begin{exampleblock}{Avantages}
        \begin{itemize}
          \item Base SQL écrite en Java, démarre en quelques millisecondes, dans le même processus
          \item Rien à installer, pas de Docker
          \item Utilisée par défaut par Spring dans un \texttt{@DataJpaTest} si elle est présente
        \end{itemize}
      \end{exampleblock}
    \end{column}
    \begin{column}{0.48\textwidth}
      \begin{alertblock}{Mais ce n'est pas PostgreSQL}
        \begin{itemize}
          \item Pas de PL/pgSQL : vos procédures stockées et triggers de R3.07 ne s'exécutent pas
          \item Types et fonctions propres à PostgreSQL absents ou imités (\texttt{jsonb}, fonctions de date\ldots)
          \item Des comportements \textbf{différents sans erreur} : tri des \texttt{NULL}, comparaisons de chaînes\ldots
          \item Un « mode de compatibilité PostgreSQL » existe, mais reste une imitation
        \end{itemize}
      \end{alertblock}
    \end{column}
  \end{columns}
  \vspace{3mm}
  \begin{center}
    Un test vert sur H2 prouve que le code marche\ldots{} \textbf{sur H2}.
  \end{center}
\end{frame}

\begin{frame}[fragile]{Un exemple : où vont les \texttt{NULL} ?}
  \begin{columns}[T]
    \begin{column}{0.55\textwidth}
\begin{lstlisting}[basicstyle=\lstcode\scriptsize]
public interface CompteRepository
        extends JpaRepository<Compte, Long> {
    // plafondRetrait == null : pas de plafond
    List<Compte> findAllByOrderByPlafondRetraitAsc();
}
\end{lstlisting}
\begin{lstlisting}[basicstyle=\lstcode\scriptsize]
comptes.save(new Compte("FR76-001", ZERO, eur("500")));
comptes.save(new Compte("FR76-002", ZERO, null));
comptes.save(new Compte("FR76-003", ZERO, eur("100")));

assertThat(comptes.findAllByOrderByPlafondRetraitAsc())
    .extracting(Compte::getNumero)
    .containsExactly("FR76-003", "FR76-001", "FR76-002");
\end{lstlisting}
    \end{column}
    \begin{column}{0.42\textwidth}
      \begin{datacard}
        \small
        \begin{tabular}{@{}ll@{}}
          \textbf{PostgreSQL} & 003, 001, \alert{002} \\
          \textbf{H2} & \alert{002}, 003, 001 \\
        \end{tabular}
      \end{datacard}
      \begin{itemize}
        \item Même code, même requête générée, \textbf{aucune erreur}
        \item PostgreSQL range les \texttt{NULL} à la fin d'un tri croissant ; H2 au début
        \item Un écran « comptes par plafond » serait faux en production alors que tous les tests passent
      \end{itemize}
    \end{column}
  \end{columns}
\end{frame}

\begin{frame}{Option 2 : Testcontainers, une vraie base jetable}
  \begin{columns}[c]
    \begin{column}{0.52\textwidth}
      \begin{center}
      \begin{tikzpicture}[node distance=0.45cm]
        \node[comp] (test) {Test JUnit};
        \node[dep, below=of test] (tc) {Testcontainers};
        \node[dep, below=of tc] (docker) {Docker};
        \node[comp, right=1.2cm of docker, draw=qaaccent2] (pg) {conteneur\\\texttt{postgres:18}};
        \draw[arr] (test) -- node[right, font=\tiny, text=qamuted] {1. démarre} (tc);
        \draw[arr] (tc) -- (docker);
        \draw[arr] (docker) -- node[above, font=\tiny, text=qamuted] {2. lance} (pg);
        \draw[arr, qaaccent2] (test.east) -| node[pos=0.3, above, font=\tiny, text=qamuted] {3. JDBC (port aléatoire)} (pg.north);
      \end{tikzpicture}
      \end{center}
    \end{column}
    \begin{column}{0.46\textwidth}
      \begin{itemize}
        \item Bibliothèque Java qui démarre des conteneurs Docker \textbf{depuis le test}
        \item La \textbf{même} version de PostgreSQL qu'en production
        \item Neuve à chaque exécution, détruite à la fin : aucune donnée qui traîne
        \item Port choisi au hasard : pas de conflit avec votre base de développement
        \item Existe pour MySQL, MongoDB, Kafka, Redis, un faux serveur HTTP\ldots
        \item Coût : Docker obligatoire, quelques secondes au démarrage
      \end{itemize}
    \end{column}
  \end{columns}
\end{frame}

\begin{frame}[fragile]{Testcontainers avec Spring Boot}
  \begin{columns}[T]
    \begin{column}{0.62\textwidth}
\begin{lstlisting}[basicstyle=\lstcode\scriptsize]
@DataJpaTest
@AutoConfigureTestDatabase(replace = Replace.NONE)
@Testcontainers
class CompteRepositoryPostgresTest {
  @Container @ServiceConnection
  static PostgreSQLContainer postgres =
      new PostgreSQLContainer("postgres:18-alpine");

  @Autowired CompteRepository comptes;

  @Test
  void comptesADecouvert_requeteDerivee() {
    comptes.save(new Compte("FR76-003", eur("-15.00")));
    comptes.save(new Compte("FR76-002", eur("0.00")));

    assertThat(comptes.findBySoldeLessThanOrderByNumero(ZERO))
        .extracting(Compte::getNumero)
        .containsExactly("FR76-003");
  }
}
\end{lstlisting}
    \end{column}
    \begin{column}{0.35\textwidth}
      \small
      \begin{itemize}
        \item \texttt{@DataJpaTest} : seulement JPA et les repositories
        \item \texttt{Replace.NONE} : ne pas remplacer la base par H2
        \item \texttt{@Testcontainers} + \texttt{@Container} : démarrer et arrêter le conteneur
        \item \texttt{@ServiceConnection} : Spring lit l'URL, l'utilisateur et le mot de passe du conteneur
        \item \texttt{static} : un conteneur pour toute la classe
      \end{itemize}
    \end{column}
  \end{columns}
\end{frame}

\begin{frame}{Isoler les tests les uns des autres}
  \begin{columns}[T]
    \begin{column}{0.48\textwidth}
      \begin{block}{Le problème}
        \begin{itemize}
          \item Le test A crée le compte \texttt{FR76-001}
          \item Le test B crée aussi \texttt{FR76-001}\ldots{} \alert{violation de contrainte \texttt{UNIQUE}}
          \item B passe seul, échoue après A : le résultat dépend de l'\textbf{ordre} (le I de FIRST)
        \end{itemize}
      \end{block}
    \end{column}
    \begin{column}{0.48\textwidth}
      \begin{block}{Les solutions}
        \begin{itemize}
          \item \textbf{Rollback} : chaque test s'exécute dans une transaction annulée à la fin (par défaut avec \texttt{@DataJpaTest}, ou \texttt{@Transactional} sur la classe de test)
          \item \textbf{Nettoyage} explicite (\texttt{deleteAll()} dans un \texttt{@BeforeEach}) quand le code testé gère lui-même ses transactions
          \item Chaque test crée \textbf{ses} données : ne jamais compter sur celles d'un autre test
        \end{itemize}
      \end{block}
    \end{column}
  \end{columns}
\end{frame}

\begin{frame}[fragile]{Des jeux de données reproductibles}
  \begin{columns}[T]
    \begin{column}{0.55\textwidth}
\begin{lstlisting}[basicstyle=\lstcode\scriptsize]
@Test
@Sql("/donnees/trois-comptes.sql")
void virement_entreDeuxComptesExistants() { ... }
\end{lstlisting}
\begin{lstlisting}[language=SQL, basicstyle=\lstcode\scriptsize]
-- src/test/resources/donnees/trois-comptes.sql
INSERT INTO compte (id, numero, solde)
VALUES (1, 'FR76-001', 500.00),
       (2, 'FR76-002', 0.00),
       (3, 'FR76-003', -15.00);
\end{lstlisting}
    \end{column}
    \begin{column}{0.42\textwidth}
      \begin{itemize}
        \item \texttt{@Sql} : exécute un script avant le test (annulé avec lui par le rollback)
        \item Ou créer les données en Java, avec des méthodes utilitaires (\texttt{unCompte().avecSolde("500")})
        \item Données \textbf{minimales} : juste ce dont le test a besoin, lisibles dans le test
        \item \alert{Jamais} de copie des données de production (RGPD, et elles changent)
      \end{itemize}
    \end{column}
  \end{columns}
\end{frame}

% ================================================================
\section{Tests d'intégration avec Spring}

\begin{frame}{Charger plus ou moins d'application}
  \footnotesize
  \begin{datacard}
    \begin{tabular}{@{}p{3.0cm}p{4.6cm}p{1.8cm}p{2.6cm}@{}}
      \toprule
      \textbf{Annotation} & \textbf{Ce qui est chargé} & \textbf{Durée} & \textbf{Pour tester} \\
      \midrule
      \textit{aucune} (JUnit seul) & rien de Spring & ms & domaine, services avec mocks \\[0.5mm]
      \texttt{@DataJpaTest} & JPA, repositories, base de test & secondes & requêtes, mappings, contraintes \\[0.5mm]
      \texttt{@WebMvcTest} & contrôleurs, JSON (services à simuler) & secondes & URL, statuts HTTP, JSON \\[0.5mm]
      \texttt{@SpringBootTest} & \textbf{toute} l'application & plusieurs secondes & parcours complet \\
      \bottomrule
    \end{tabular}
  \end{datacard}
  \vspace{1mm}
  \begin{itemize}\small
    \item Les « tranches » (\textit{slices}) \texttt{@DataJpaTest} et \texttt{@WebMvcTest} ne chargent qu'une couche : plus rapides, échecs plus faciles à localiser
    \item Spring \textbf{réutilise} le contexte entre classes de test qui ont la même configuration : un \texttt{@MockitoBean} différent dans chaque classe oblige à le recréer
  \end{itemize}
\end{frame}

\begin{frame}[fragile]{Toute l'application, sauf l'API externe}
\begin{lstlisting}[basicstyle=\lstcode\scriptsize]
@SpringBootTest @AutoConfigureMockMvc @Testcontainers @Transactional
class CompteControllerIntegrationTest {
    @Container @ServiceConnection
    static PostgreSQLContainer postgres = new PostgreSQLContainer("postgres:18-alpine");

    @Autowired MockMvc mvc;                   // envoie de fausses requêtes HTTP, sans serveur
    @Autowired CompteRepository comptes;
    @MockitoBean TauxChangeClient taux;       // remplace le vrai bean dans le contexte Spring
    @Test
    void soldeEnDollars() throws Exception {
        comptes.save(new Compte("FR76-001", eur("100.00")));
        when(taux.taux("EUR", "USD")).thenReturn(eur("1.10"));
        mvc.perform(get("/api/comptes/FR76-001/solde").param("devise", "USD"))
           .andExpect(status().isOk())
           .andExpect(jsonPath("$.solde").value(110.00))
           .andExpect(jsonPath("$.devise").value("USD"));
    }
}
\end{lstlisting}
  \source{\texttt{@MockitoBean} remplace l'ancien \texttt{@MockBean}. Contrôleur, service, repository, JPA et PostgreSQL sont réels.}
\end{frame}

\begin{frame}[fragile]{Une configuration par environnement}
  \begin{columns}[T]
    \begin{column}{0.5\textwidth}
\begin{lstlisting}[language={}, basicstyle=\lstcode\scriptsize]
# src/main/resources/application.properties
spring.datasource.url=jdbc:postgresql://localhost:5432/banque
spring.datasource.username=banque

# src/test/resources/application-test.properties
spring.jpa.show-sql=true
banque.notifications.actives=false
\end{lstlisting}
\begin{lstlisting}[basicstyle=\lstcode\scriptsize]
@SpringBootTest
@ActiveProfiles("test")
class ... { }
\end{lstlisting}
    \end{column}
    \begin{column}{0.46\textwidth}
      \small
      \begin{itemize}
        \item Le \textbf{code} est le même partout ; la \textbf{configuration} change selon l'environnement : développement, test, recette, production (CM6)
        \item Profils Spring : \texttt{application-<profil>.properties}, activé par \texttt{@ActiveProfiles} ou une variable d'environnement
        \item Avec \texttt{@ServiceConnection}, l'adresse de la base de test vient du conteneur : impossible de toucher la base de développement par erreur
        \item Les mots de passe de production ne sont \alert{jamais} dans le dépôt
      \end{itemize}
    \end{column}
  \end{columns}
\end{frame}

\begin{frame}[fragile]{Et l'API externe ? Tester l'adaptateur}
  \begin{columns}[T]
    \begin{column}{0.62\textwidth}
\begin{lstlisting}[basicstyle=\lstcode\scriptsize]
@WireMockTest      // faux serveur HTTP local
class HttpTauxChangeClientTest {
  @Test
  void lit_le_taux(WireMockRuntimeInfo wm) {
    stubFor(get("/v1/EUR/USD")
        .willReturn(okJson("{\"taux\": 1.10}")));
    var client = new HttpTauxChangeClient(wm.getHttpBaseUrl());
    assertThat(client.taux(EUR, USD)).isEqualByComparingTo("1.10");
  }
  @Test
  void erreur500_devientTauxIndisponible(WireMockRuntimeInfo wm) {
    stubFor(get("/v1/EUR/USD").willReturn(serverError()));
    // ... assertThatThrownBy(...) : TauxIndisponibleException
  }
}
\end{lstlisting}
    \end{column}
    \begin{column}{0.35\textwidth}
      \small
      \begin{itemize}
        \item Au CM3, on a simulé \textbf{notre} interface ; reste à tester son implémentation HTTP
        \item Pas la vraie API (payante, lente, instable) : un \textbf{faux serveur HTTP} local (WireMock)
        \item Vraies requêtes, vrai JSON, vraies erreurs (500, délai dépassé)
        \item Les réponses simulées suivent la documentation du fournisseur : à revérifier de temps en temps
      \end{itemize}
    \end{column}
  \end{columns}
\end{frame}

\rdemo{1}{Une vraie base dans le test}{%
  \item Lancer le test PostgreSQL : un conteneur apparaît (\texttt{docker ps}), le test passe, le conteneur disparaît
  \item Lancer \textbf{le même test} sur H2 : rouge, sans aucune erreur SQL\ldots{} les \texttt{NULL} ne sont pas au même endroit
  \item Toute l'application, contrôleur compris, avec seulement l'API de change simulée
}{\texttt{CompteRepositoryPostgresTest}\\\texttt{CompteRepositoryH2Test}\\\texttt{CompteControllerIntegrationTest}\\[1mm]terminal : \texttt{docker ps}}

\begin{frame}{Le coût des tests d'intégration}
  \begin{columns}[T]
    \begin{column}{0.48\textwidth}
      \begin{itemize}
        \item Plus lents : démarrage de Spring, du conteneur, accès disque et réseau
        \item Plus difficiles à diagnostiquer : un échec peut venir de n'importe quelle couche
        \item Donc \textbf{moins nombreux} : un test par requête, mapping ou parcours important, pas un par règle métier
      \end{itemize}
    \end{column}
    \begin{column}{0.48\textwidth}
      \begin{block}{Bonnes pratiques}
        \begin{itemize}
          \item Les règles métier restent testées en \textbf{unitaire} (TP1, TP3)
          \item Un conteneur \texttt{static} partagé par classe, voire par toute la suite
          \item Séparer les tests lents des tests rapides : \texttt{mvn test} pour les unitaires, \texttt{mvn verify} pour l'intégration (CM6)
          \item Les lancer \textbf{tous} en CI, à chaque modification
        \end{itemize}
      \end{block}
    \end{column}
  \end{columns}
\end{frame}

\begin{frame}{À retenir}
  \small
  \begin{itemize}
    \item Un test d'intégration vérifie que des composants \textbf{réels} fonctionnent \textbf{ensemble} : requêtes, mappings, contraintes, transactions, configuration
    \item Jamais sur la base de développement : une base \textbf{dédiée}, dans un état connu au début de chaque test
    \item H2 est rapide mais \textbf{n'est pas PostgreSQL} : différences silencieuses, pas de PL/pgSQL
    \item \textbf{Testcontainers} : la vraie base, dans Docker, jetable ; \texttt{@ServiceConnection} la branche à Spring
    \item Isolation : \textbf{rollback} après chaque test, chaque test crée ses propres données (\texttt{@Sql} ou Java)
    \item Spring : \texttt{@DataJpaTest} pour les repositories, \texttt{@SpringBootTest} + \texttt{MockMvc} pour toute l'application, \texttt{@MockitoBean} pour ce qui reste hors périmètre
    \item Une configuration par environnement (profils) ; le code ne change pas
  \end{itemize}
  \vspace{1mm}
  \begin{exampleblock}{TP4}
    Tester les repositories et l'API REST du projet fil rouge avec un vrai PostgreSQL.
  \end{exampleblock}
\end{frame}

\rquestionslide

\end{document}
